
\documentclass{article}
\usepackage{colortbl}
\usepackage{makecell}
\usepackage{multirow}
\usepackage{supertabular}

\begin{document}

\newcounter{utterance}

\centering \large Interaction Transcript for game `hot\_air\_balloon', experiment `air\_balloon\_survival\_en\_negotiation\_easy', episode 1 with mistral{-}3{-}large{-}2512{-}t1.0.
\vspace{24pt}

{ \footnotesize  \setcounter{utterance}{1}
\setlength{\tabcolsep}{0pt}
\begin{supertabular}{c@{$\;$}|p{.15\linewidth}@{}p{.15\linewidth}p{.15\linewidth}p{.15\linewidth}p{.15\linewidth}p{.15\linewidth}}
    \# & $\;$A & \multicolumn{4}{c}{Game Master} & $\;\:$B\\
    \hline

    \theutterance \stepcounter{utterance}  
    & & \multicolumn{4}{p{0.6\linewidth}}{
        \cellcolor[rgb]{0.9,0.9,0.9}{
            \makecell[{{p{\linewidth}}}]{
                \texttt{\tiny{[P1$\langle$GM]}}
                \texttt{You are participating in a collaborative negotiation game.} \\
\\ 
\texttt{Together with another participant, you must agree on a single set of items that will be kept. Each of you has your own view of how much each item matters to you (importance). You do not know how the other participant values the items. Additionally, you are given the effort each item demands.} \\
\texttt{You may only agree on a set if the total effort of the selected items does not exceed a shared limit:} \\
\\ 
\texttt{LIMIT = 4564} \\
\\ 
\texttt{Here are the individual item effort values:} \\
\\ 
\texttt{Item effort = \{"C08": 299, "C55": 241, "A43": 889, "A67": 967, "B42": 619, "A62": 799, "B19": 978, "C43": 733, "C73": 909, "C10": 501, "B36": 139, "A18": 594, "C61": 565, "C67": 789, "C42": 107\}} \\
\\ 
\texttt{Here is your personal view on the importance of each item:} \\
\\ 
\texttt{Item importance values = \{"C08": 138, "C55": 583, "A43": 868, "A67": 822, "B42": 783, "A62": 65, "B19": 262, "C43": 121, "C73": 508, "C10": 780, "B36": 461, "A18": 484, "C61": 668, "C67": 389, "C42": 808\}} \\
\\ 
\texttt{Goal:} \\
\\ 
\texttt{Your goal is to negotiate a shared set of items that benefits you as much as possible (i.e., maximizes total importance to YOU), while staying within the LIMIT. You are not required to make a PROPOSAL in every message {-} you can simply negotiate as well. All tactics are allowed!} \\
\\ 
\texttt{Interaction Protocol:} \\
\\ 
\texttt{You may only use the following structured formats in a message:} \\
\\ 
\texttt{PROPOSAL: \{'A', 'B', 'C', …\}} \\
\texttt{Propose keeping exactly those items.} \\
\texttt{REFUSE: \{'A', 'B', 'C', …\}} \\
\texttt{Explicitly reject opponent's proposal.} \\
\texttt{ARGUMENT: \{'...'\}} \\
\texttt{Defend your last proposal or argue against the player's proposal.} \\
\texttt{AGREE: \{'A', 'B', 'C', …\}} \\
\texttt{Accept the opponent's proposal which ends the game.} \\
\texttt{STRATEGIC REASONING: \{'...'\}} \\
\texttt{	Describe your strategic reasoning or anticipation explaining your choice of action. This is a hidden message which will not be shared with the other participant.} \\
\\ 
\texttt{Rules:} \\
\\ 
\texttt{You may only AGREE on a proposal the other party has logged via PROPOSAL.} \\
\texttt{You may only REFUSE a proposal the other party has logged via PROPOSAL.} \\
\texttt{Total effort of any PROPOSAL or AGREE set must be ≤ LIMIT.} \\
\texttt{Do NOT reveal your hidden importance scores.} \\
\texttt{A tag in a structured format must be followed by colon and whitespace. The argument must be a python set containing 0 or more strings.} \\
\texttt{So, it must be of the form TAG: \{...\}} \\
\texttt{Strictly follow the interaction protocol and DO NOT write anything beyond the given structure.} \\
\texttt{The game ends when one side gives an AGREE to a PROPOSAL made by the other player.} \\
\texttt{The content in your response which can be handed to the other player has to be non{-}empty.} \\
\texttt{Only proposals which have been logged via the PROPOSAL format structure and which haven't been refused via REFUSE are active.} \\
\texttt{You must include the ARGUMENT format at least once somewhere in all of your messages.} \\
\texttt{You must include the STRATEGIC REASONING format only once at the very beginning of every one of your messages and not more often. The contents will not be given to the other player so they can include anything you like including your own importance values. Here you should reason step by step to come up with you next move.} \\
            }
        }
    }
    & \\ \\

    \theutterance \stepcounter{utterance}  
    & \multicolumn{4}{p{0.6\linewidth}}{
        \cellcolor[rgb]{0.9,0.9,0.9}{
            \makecell[{{p{\linewidth}}}]{
                \texttt{\tiny{[P1$\rangle$GM]}}
                \texttt{STRATEGIC REASONING: \{} \\
\texttt{I need to maximize my total importance while staying under the LIMIT of 4564. First, I should identify the items that give me the highest importance{-}to{-}effort ratio, as these provide the most "bang for the buck." Here’s a quick ranking of items by importance/effort (higher is better):} \\
\\ 
\texttt{{-} C42: 808/107 ≈ 7.55} \\
\texttt{{-} B36: 461/139 ≈ 3.32} \\
\texttt{{-} C55: 583/241 ≈ 2.42} \\
\texttt{{-} C10: 780/501 ≈ 1.56} \\
\texttt{{-} C61: 668/565 ≈ 1.18} \\
\texttt{{-} A18: 484/594 ≈ 0.81} \\
\texttt{{-} C08: 138/299 ≈ 0.46} \\
\texttt{{-} C73: 508/909 ≈ 0.56} \\
\texttt{{-} B42: 783/619 ≈ 1.26} \\
\texttt{{-} A43: 868/889 ≈ 0.98} \\
\texttt{{-} C67: 389/789 ≈ 0.49} \\
\texttt{{-} A67: 822/967 ≈ 0.85} \\
\texttt{{-} C43: 121/733 ≈ 0.17} \\
\texttt{{-} A62: 65/799 ≈ 0.08} \\
\texttt{{-} B19: 262/978 ≈ 0.27} \\
\\ 
\texttt{The top items by this ratio are C42, B36, C55, and C10. I should prioritize these in negotiations, as they give me the most value per unit of effort. However, I must also consider that the other player’s preferences are unknown, so I should avoid revealing my true importance values.} \\
\\ 
\texttt{For my initial proposal, I will include the highest ratio items first, then fill the remaining effort with other high{-}importance items. The total effort of C42 (107) + B36 (139) + C55 (241)} \\
            }
        }
    }
    & & \\ \\

    \theutterance \stepcounter{utterance}  
    & & \multicolumn{4}{p{0.6\linewidth}}{
        \cellcolor[rgb]{0.9,0.9,0.9}{
            \makecell[{{p{\linewidth}}}]{
                \texttt{\tiny{[P1$\langle$GM]}}
                \texttt{Your response did not start with the proper strategic reasoning tag at the very beginning of your response. The very first structured format must be of the form STRATEGIC REASONING: \{...\}. Try again.} \\
            }
        }
    }
    & \\ \\

    \theutterance \stepcounter{utterance}  
    & \multicolumn{4}{p{0.6\linewidth}}{
        \cellcolor[rgb]{0.9,0.9,0.9}{
            \makecell[{{p{\linewidth}}}]{
                \texttt{\tiny{[P1$\rangle$GM]}}
                \texttt{STRATEGIC REASONING: \{} \\
\texttt{My goal is to maximize my total importance while staying under the LIMIT of 4564. I will prioritize items with the highest importance{-}to{-}effort ratio, as these give me the most value per unit of effort. Here’s my plan:} \\
\\ 
\texttt{1. **Rank items by importance/effort ratio** (descending):} \\
\texttt{   {-} C42: 808/107 ≈ 7.55} \\
\texttt{   {-} B36: 461/139 ≈ 3.32} \\
\texttt{   {-} C55: 583/241 ≈ 2.42} \\
\texttt{   {-} C10: 780/501 ≈ 1.56} \\
\texttt{   {-} B42: 783/619 ≈ 1.26} \\
\texttt{   {-} C61: 668/565 ≈ 1.18} \\
\texttt{   {-} A18: 484/594 ≈ 0.81} \\
\texttt{   {-} A67: 822/967 ≈ 0.85} \\
\texttt{   {-} A43: 868/889 ≈ 0.98} \\
\texttt{   {-} C73: 508/909 ≈ 0.56} \\
\texttt{   {-} C08: 138/299 ≈ 0.46} \\
\texttt{   {-} C67: 389/789 ≈ 0.49} \\
\texttt{   {-} B19: 262/978 ≈ 0.27} \\
\texttt{   {-} C43: 121/733 ≈ 0.17} \\
\texttt{   {-} A62: 65/799 ≈ 0.08} \\
\\ 
\texttt{2. **Initial proposal**: Start with the highest{-}ratio items (C42, B36, C55, C10) and fill remaining effort with other high{-}importance items (B42, C61, A18). This gives:} \\
\texttt{   {-} C42 (107) + B36 (139) + C55 (241) + C10 (501) + B42 (619) + C61 (56} \\
            }
        }
    }
    & & \\ \\

    \theutterance \stepcounter{utterance}  
    & & \multicolumn{4}{p{0.6\linewidth}}{
        \cellcolor[rgb]{0.9,0.9,0.9}{
            \makecell[{{p{\linewidth}}}]{
                \texttt{\tiny{[P1$\langle$GM]}}
                \texttt{Your response did not start with the proper strategic reasoning tag at the very beginning of your response. The very first structured format must be of the form STRATEGIC REASONING: \{...\}. Try again.} \\
            }
        }
    }
    & \\ \\

    \theutterance \stepcounter{utterance}  
    & \multicolumn{4}{p{0.6\linewidth}}{
        \cellcolor[rgb]{0.9,0.9,0.9}{
            \makecell[{{p{\linewidth}}}]{
                \texttt{\tiny{[P1$\rangle$GM]}}
                \texttt{STRATEGIC REASONING: \{} \\
\texttt{**Step 1: Calculate optimal set for maximum personal importance**} \\
\texttt{{-} Sort items by importance/effort ratio (descending) to prioritize high{-}value, low{-}effort items.} \\
\texttt{{-} Greedy selection: Start with C42 (7.55), B36 (3.32), C55 (2.42), C10 (1.56), B42 (1.26), C61 (1.18), A18 (0.81).} \\
\texttt{{-} Total effort: 107+139+241+501+619+565+594 = 2766 (well under LIMIT).} \\
\texttt{{-} Total importance: 808+461+583+780+783+668+484 = 4567 (very high for me).} \\
\\ 
\texttt{**Step 2: Anticipate opponent’s behavior**} \\
\texttt{{-} They likely prioritize their own high{-}importance items, which may overlap with mine (e.g., A43, A67, B19).} \\
\texttt{{-} If they propose excluding my top items (C42, B36), I must argue for their inclusion by framing them as "low{-}effort, high{-}value" without revealing my true importance.} \\
\\ 
\texttt{**Step 3: Initial proposal strategy**} \\
\texttt{{-} Propose the optimal set (C42, B36, C55, C10, B42, C61, A18) to anchor negotiations.} \\
\texttt{{-} If rejected, use ARGUMENT to emphasize "efficiency" and "shared benefit" of these items.} \\
\texttt{{-} If they counter with exclusions, prioritize keeping C42, B36, C55, and C10 at all costs.} \\
\texttt{\}} \\
\\ 
\texttt{PROPOSAL: \{'C42', 'B36', 'C55', 'C10', 'B42', 'C61', 'A18'\}} \\
\texttt{ARGUMENT: \{} \\
\texttt{These items maximize efficiency (low effort, high collective value) and leave ample room under the LIMIT for adjustments. Their combined effort (2766) is far below the threshold, allowing flexibility for additional items if needed.} \\
\texttt{\}} \\
            }
        }
    }
    & & \\ \\

    \theutterance \stepcounter{utterance}  
    & & \multicolumn{4}{p{0.6\linewidth}}{
        \cellcolor[rgb]{0.9,0.9,0.9}{
            \makecell[{{p{\linewidth}}}]{
                \texttt{\tiny{[GM$\rangle$P2]}}
                \texttt{You are participating in a collaborative negotiation game.} \\
\\ 
\texttt{Together with another participant, you must agree on a single set of items that will be kept. Each of you has your own view of how much each item matters to you (importance). You do not know how the other participant values the items. Additionally, you are given the effort each item demands.} \\
\texttt{You may only agree on a set if the total effort of the selected items does not exceed a shared limit:} \\
\\ 
\texttt{LIMIT = 4564} \\
\\ 
\texttt{Here are the individual item effort values:} \\
\\ 
\texttt{Item effort = \{"C08": 299, "C55": 241, "A43": 889, "A67": 967, "B42": 619, "A62": 799, "B19": 978, "C43": 733, "C73": 909, "C10": 501, "B36": 139, "A18": 594, "C61": 565, "C67": 789, "C42": 107\}} \\
\\ 
\texttt{Here is your personal view on the importance of each item:} \\
\\ 
\texttt{Item importance values = \{"C08": 138, "C55": 583, "A43": 868, "A67": 822, "B42": 783, "A62": 65, "B19": 262, "C43": 121, "C73": 508, "C10": 780, "B36": 461, "A18": 484, "C61": 668, "C67": 389, "C42": 808\}} \\
\\ 
\texttt{Goal:} \\
\\ 
\texttt{Your goal is to negotiate a shared set of items that benefits you as much as possible (i.e., maximizes total importance to YOU), while staying within the LIMIT. You are not required to make a PROPOSAL in every message {-} you can simply negotiate as well. All tactics are allowed!} \\
\\ 
\texttt{Interaction Protocol:} \\
\\ 
\texttt{You may only use the following structured formats in a message:} \\
\\ 
\texttt{PROPOSAL: \{'A', 'B', 'C', …\}} \\
\texttt{Propose keeping exactly those items.} \\
\texttt{REFUSE: \{'A', 'B', 'C', …\}} \\
\texttt{Explicitly reject opponent's proposal.} \\
\texttt{ARGUMENT: \{'...'\}} \\
\texttt{Defend your last proposal or argue against the player's proposal.} \\
\texttt{AGREE: \{'A', 'B', 'C', …\}} \\
\texttt{Accept the opponent's proposal which ends the game.} \\
\texttt{STRATEGIC REASONING: \{'...'\}} \\
\texttt{	Describe your strategic reasoning or anticipation explaining your choice of action. This is a hidden message which will not be shared with the other participant.} \\
\\ 
\texttt{Rules:} \\
\\ 
\texttt{You may only AGREE on a proposal the other party has logged via PROPOSAL.} \\
\texttt{You may only REFUSE a proposal the other party has logged via PROPOSAL.} \\
\texttt{Total effort of any PROPOSAL or AGREE set must be ≤ LIMIT.} \\
\texttt{Do NOT reveal your hidden importance scores.} \\
\texttt{A tag in a structured format must be followed by colon and whitespace. The argument must be a python set containing 0 or more strings.} \\
\texttt{So, it must be of the form TAG: \{...\}} \\
\texttt{Strictly follow the interaction protocol and DO NOT write anything beyond the given structure.} \\
\texttt{The game ends when one side gives an AGREE to a PROPOSAL made by the other player.} \\
\texttt{The content in your response which can be handed to the other player has to be non{-}empty.} \\
\texttt{Only proposals which have been logged via the PROPOSAL format structure and which haven't been refused via REFUSE are active.} \\
\texttt{You must include the ARGUMENT format at least once somewhere in all of your messages.} \\
\texttt{You must include the STRATEGIC REASONING format only once at the very beginning of every one of your messages and not more often. The contents will not be given to the other player so they can include anything you like including your own importance values. Here you should reason step by step to come up with you next move.} \\
\\ 
\texttt{You will now receive the first message of the other player.} \\
\\ 
\texttt{PROPOSAL: \{'C42', 'B36', 'C55', 'C10', 'B42', 'C61', 'A18'\}} \\
\texttt{ARGUMENT: \{} \\
\texttt{These items maximize efficiency (low effort, high collective value) and leave ample room under the LIMIT for adjustments. Their combined effort (2766) is far below the threshold, allowing flexibility for additional items if needed.} \\
\texttt{\}} \\
            }
        }
    }
    & \\ \\

    \theutterance \stepcounter{utterance}  
    & & & \multicolumn{4}{p{0.6\linewidth}}{
        \cellcolor[rgb]{0.9,0.9,0.9}{
            \makecell[{{p{\linewidth}}}]{
                \texttt{\tiny{[GM$\langle$P2]}}
                \texttt{STRATEGIC REASONING: \{} \\
\texttt{The opponent's proposal includes items with a total effort of 2766 (C42:107 + B36:139 + C55:241 + C10:501 + B42:619 + C61:565 + A18:594). While this is under the LIMIT, I need to evaluate how well this aligns with my importance scores.} \\
\\ 
\texttt{Key observations:} \\
\texttt{1. My high{-}importance items not included: C08 (138), A43 (868), A67 (822), C73 (508), C67 (389).} \\
\texttt{2. Included items with low importance to me: A18 (484), C61 (668), B42 (783) are decent, but C42 (808) is actually high for me.} \\
\texttt{3. The proposal leaves 1798 effort unused (4564 {-} 2766). This could accommodate more of my high{-}importance items.} \\
\\ 
\texttt{Tactics:} \\
\texttt{{-} I should counter{-}propose by replacing some of their lower{-}importance{-}to{-}me items with my high{-}importance items while staying under the LIMIT.} \\
\texttt{{-} Specifically, I can remove A18 (594 effort, 484 importance) and C61 (565 effort, 668 importance) to free up 1159 effort, then add A43 (889 effort, 868 importance) and A67 (967 effort, 822 importance). This would exceed the LIMIT, so I'll need to adjust.} \\
\texttt{{-} A better approach: Remove A18 (594) and add A43 (889) and C08 (299). Total effort: 2766 {-} 594 + 889 + 299 = 3360 (under LIMIT). This adds 868 + 138 = 1006 importance vs. 484 lost, a net gain of 522.} \\
\texttt{{-} Alternatively, remove C61 (565) and add A67 (967). Total effort: 2766} \\
            }
        }
    }
     \\ \\

    \theutterance \stepcounter{utterance}  
    & & \multicolumn{4}{p{0.6\linewidth}}{
        \cellcolor[rgb]{0.9,0.9,0.9}{
            \makecell[{{p{\linewidth}}}]{
                \texttt{\tiny{[GM$\rangle$P2]}}
                \texttt{Your response did not start with the proper strategic reasoning tag at the very beginning of your response. The very first structured format must be of the form STRATEGIC REASONING: \{...\}. Try again.} \\
            }
        }
    }
    & \\ \\

    \theutterance \stepcounter{utterance}  
    & & & \multicolumn{4}{p{0.6\linewidth}}{
        \cellcolor[rgb]{0.9,0.9,0.9}{
            \makecell[{{p{\linewidth}}}]{
                \texttt{\tiny{[GM$\langle$P2]}}
                \texttt{STRATEGIC REASONING: \{} \\
\texttt{Re{-}evaluating the opponent's proposal:} \\
\texttt{{-} Their proposal: \{'C42', 'B36', 'C55', 'C10', 'B42', 'C61', 'A18'\} (effort = 2766, my importance = 808+461+583+780+783+668+484 = 4567).} \\
\texttt{{-} My high{-}importance items missing: A43 (868), A67 (822), C73 (508), C08 (138), C67 (389).} \\
\texttt{{-} The proposal leaves 1798 effort unused (4564 {-} 2766). I can use this to add my high{-}importance items.} \\
\\ 
\texttt{Optimal counter{-}proposal:} \\
\texttt{1. Remove A18 (594 effort, 484 importance) and C61 (565 effort, 668 importance) to free up 1159 effort.} \\
\texttt{2. Add A43 (889 effort, 868 importance) and C08 (299 effort, 138 importance). New effort: 2766 {-} 594 {-} 565 + 889 + 299 = 2795 (under LIMIT).} \\
\texttt{   {-} Net importance gain: (868 + 138) {-} (484 + 668) = {-}146 (bad).} \\
\texttt{3. Better: Remove A18 (594) and add A43 (889) and C73 (909). New effort: 2766 {-} 594 + 889 + 909 = 3970 (under LIMIT).} \\
\texttt{   {-} Net importance gain: (868 + 508) {-} 484 = 892 (good).} \\
\texttt{4. Even better: Remove C61 (565) and add A67 (967). New effort: 2766 {-} 565 + 967 = 3168 (under LIMIT).} \\
\texttt{   {-} Net importance gain: 82} \\
            }
        }
    }
     \\ \\

    \theutterance \stepcounter{utterance}  
    & & \multicolumn{4}{p{0.6\linewidth}}{
        \cellcolor[rgb]{0.9,0.9,0.9}{
            \makecell[{{p{\linewidth}}}]{
                \texttt{\tiny{[GM$\rangle$P2]}}
                \texttt{Your response did not start with the proper strategic reasoning tag at the very beginning of your response. The very first structured format must be of the form STRATEGIC REASONING: \{...\}. Try again.} \\
            }
        }
    }
    & \\ \\

    \theutterance \stepcounter{utterance}  
    & & & \multicolumn{4}{p{0.6\linewidth}}{
        \cellcolor[rgb]{0.9,0.9,0.9}{
            \makecell[{{p{\linewidth}}}]{
                \texttt{\tiny{[GM$\langle$P2]}}
                \texttt{STRATEGIC REASONING: \{} \\
\texttt{Let's analyze the opponent's proposal and plan a counter{-}strategy:} \\
\\ 
\texttt{1. Opponent's proposal: \{'C42', 'B36', 'C55', 'C10', 'B42', 'C61', 'A18'\}} \\
\texttt{   {-} Total effort: 107+139+241+501+619+565+594 = 2766} \\
\texttt{   {-} My importance: 808+461+583+780+783+668+484 = 4567} \\
\texttt{   {-} Remaining effort: 4564{-}2766 = 1798} \\
\\ 
\texttt{2. My high{-}value items missing from their proposal:} \\
\texttt{   {-} A43: 889 effort, 868 importance} \\
\texttt{   {-} A67: 967 effort, 822 importance} \\
\texttt{   {-} C73: 909 effort, 508 importance} \\
\texttt{   {-} C08: 299 effort, 138 importance} \\
\\ 
\texttt{3. Best counter{-}proposal approach:} \\
\texttt{   {-} Remove A18 (594 effort, 484 importance) and C61 (565 effort, 668 importance)} \\
\texttt{   {-} Add A43 (889 effort, 868 importance) and A67 (967 effort, 822 importance)} \\
\texttt{   {-} New effort: 2766 {-} 594 {-} 565 + 889 + 967 = 3463 (under LIMIT)} \\
\texttt{   {-} Net importance gain: (868+822) {-} (484+668) = 538} \\
\\ 
\texttt{4. Alternative if they resist:} \\
\texttt{   {-} Remove just A18 (594 effort) and add A43 (889 effort) and C08 (299 effort)} \\
\texttt{   {-} New effort: 2766 {-} 594 + 889 + 299 = 3360} \\
\texttt{   {-} Net importance gain: (868+138) {-} 484 = 522} \\
\\ 
\texttt{5. Negoti} \\
            }
        }
    }
     \\ \\

    \theutterance \stepcounter{utterance}  
    & & & \multicolumn{2}{p{0.3\linewidth}}{
        \cellcolor[rgb]{0.9,0.9,0.9}{
            \makecell[{{p{\linewidth}}}]{
                \texttt{\tiny{[GM$|$GM]}}
                \texttt{abort} \\
            }
        }
    }
    & & \\ \\

    \theutterance \stepcounter{utterance}  
    & & & \multicolumn{2}{p{0.3\linewidth}}{
        \cellcolor[rgb]{0.9,0.9,0.9}{
            \makecell[{{p{\linewidth}}}]{
                \texttt{\tiny{[GM$|$GM]}}
                \texttt{end game} \\
            }
        }
    }
    & & \\ \\

\end{supertabular}
}

\end{document}
